% ============================================================
%  SUPPLEMENT 06 — FLOATING-POINT PRECISION
%  float32 vs float64: When It Matters and When to Cast
%  Companion to: ECDR_Data_Analysis.tex
% ============================================================
\documentclass[11pt, a4paper]{article}

\usepackage{fontspec}
\setmainfont{Noto Serif}
\setsansfont{Noto Sans}
\setmonofont{Noto Sans Mono}[Scale=0.88]

\usepackage[english]{babel}
\usepackage{amsmath, amssymb, amsthm}
\usepackage{mathtools}
\usepackage[margin=2.5cm]{geometry}
\usepackage{booktabs, tabularx, array}
\usepackage{xcolor}
\newcolumntype{Y}{>{\raggedright\arraybackslash}X}
\newcolumntype{P}[1]{>{\raggedright\arraybackslash}p{#1}}
\usepackage{hyperref,xurl}
\usepackage{enumitem}
\usepackage{microtype}
\usepackage{parskip}
\setlength{\emergencystretch}{3em}

\definecolor{ink}{RGB}{20,20,20}
\definecolor{teal}{RGB}{0,120,120}
\definecolor{muted}{RGB}{90,90,90}

\hypersetup{colorlinks=true, linkcolor=teal, urlcolor=teal, citecolor=teal}

\theoremstyle{definition}
\newtheorem{definition}{Definition}[section]
\newtheorem{example}{Example}[section]
\newtheorem{remark}{Remark}[section]
\theoremstyle{plain}
\newtheorem{theorem}{Theorem}[section]

\newcommand{\R}{\mathbb{R}}
\newcommand{\code}[1]{{\small\ttfamily\nolinkurl{#1}}}
\newcommand{\eps}{\varepsilon}

\usepackage{titlesec}
\titleformat{\section}{\large\bfseries\color{teal}}{}{0em}{\thesection\quad}
\titleformat{\subsection}{\normalsize\bfseries\color{ink}}{}{0em}{\thesubsection\quad}

\begin{document}

\begin{center}
  {\LARGE\bfseries\color{teal} Floating-Point Precision in the ECDR Pipeline}\\[0.5em]
  {\large\color{muted} float32 vs.\ float64: Semantics, Rounding, and Certification Obligations}\\[1.2em]
  {\normalsize Supplement~06 to \textit{ECDR Data Analysis} \textbar{} October~2026}\\[0.3em]
  {\small\color{muted} This document is supplementary only. The main document is not modified.}
\end{center}

\vspace{1em}\hrule\vspace{1em}

\tableofcontents
\clearpage

% ============================================================
\section{IEEE 754 Floating-Point Formats}

Both float32 and float64 follow the IEEE~754 standard for binary floating-point
arithmetic. They differ in the number of bits allocated to each part of the number:

\begin{center}
\small
\begin{tabular}{@{}lrrrl@{}}
\toprule
Format & Total bits & Sign & Exponent & Mantissa (significand) \\
\midrule
float32 (single) & 32 & 1 & 8 & 23 (+1 implicit) \\
float64 (double) & 64 & 1 & 11 & 52 (+1 implicit) \\
\bottomrule
\end{tabular}
\end{center}

A floating-point number is represented as $(-1)^s \times 1.m \times 2^{e - \text{bias}}$,
where $s$ is the sign bit, $m$ is the binary fraction stored in the mantissa, and
$e$ is the stored exponent.

% ============================================================
\section{Machine Epsilon and Rounding Error}

\begin{definition}[Machine Epsilon $\eps_{\text{mach}}$]
  The machine epsilon $\eps_{\text{mach}}$ is the distance from 1.0 to the next
  representable floating-point number:
  \[
    \eps_{\text{mach}} = 2^{-\text{mantissa bits}}.
  \]
\end{definition}

\begin{center}
\small
\begin{tabular}{@{}lrl@{}}
\toprule
Format & $\eps_{\text{mach}}$ & Approximate value \\
\midrule
float32 & $2^{-23}$ & $1.19 \times 10^{-7}$ \\
float64 & $2^{-52}$ & $2.22 \times 10^{-16}$ \\
\bottomrule
\end{tabular}
\end{center}

For a finite normal exact result $x$ rounded to nearest with no overflow
or underflow, the standard model is $\mathrm{fl}(x)=x(1+\delta)$,
$|\delta|\le u=\epsilon_{\rm mach}/2$. If $x$ is already representable,
rounding it to the same format changes nothing. Subnormal handling and
backend fused operations require additional assumptions. A sequential dot
product of length$n$ admits the conditional forward-error bound
\[
 |\mathrm{fl}(a^\top b)-a^\top b|\le\gamma_n\sum_i|a_ib_i|,
 \qquad \gamma_n=\frac{nu}{1-nu},\quad nu<1,
\]
under the specified arithmetic model. It is an absolute bound scaled by the
sum of absolute products, not a small relative error of a cancellation-prone
result. A square-root-of-n random-error rule is a statistical heuristic,
not a worst-case certificate. Library convolution/reduction order must be
justified separately before substituting$n$ into this formula.

% ============================================================
\section{Cast Semantics: float32 Rounding}

\subsection{What Casting Means}

When a float64 value $x$ is \emph{cast} to float32, it is rounded to the nearest
float32 representable value:
\[
  \text{fl}_{32}(x) = \text{round}_{32}(x).
\]
For finite normal values under round-to-nearest and no overflow/underflow,
the relative error is at most$u=2^{-24}$. For the illustrative reading
$0.82\,g$, the absolute cast error is bounded by roughly$4.9\times10^{-8}\,g$
under those assumptions. This does not establish a sensor-resolution claim
or justify accepting a certificate boundary.

\subsection{The Audit's float32 Cast Check}

Historical Evidence D01 reconstructed13,403,904 float32 scalar values
from raw text and reported zero discrepancies. It compared a reconstructed
cast path with the saved tensor representation of that historical audit;
that mention of a PyTorch\code{.pt} file is historical and does not identify
the current\code{data/windows.npz} input artifact. Current Week2 receipts
in\code{output/week2_foundations_20261103/} identify the raw-slice and
normalized-overlap checks for the K=12 pipeline.

A reconstruction check reads text, applies the stated dtype conversions,
compares saved tensor values, and counts discrepancies. Equal final values
support consistency with this reconstruction; they do not identify a unique
construction pipeline. Reading directly into float32 or using other
intermediates can produce the same values. Builder source, hashes, and
execution provenance establish which operations actually ran.

\begin{remark}[Why This Matters Semantically]
Data-value equality, artifact identity and execution provenance answer
separate questions. A count of zero mismatches does not prove the absence
of intermediate rounding or processing, and it does not prove a universal
numerical certificate. Preserve dtype and operator semantics as explicit
inputs to the analysis rather than inferring them from matching outputs.
\end{remark}

% ============================================================
\section{When Precision Matters for ECDR}

\subsection{Case 1: Distance Computation}

The distance metrics in the data audit ($d_2$, $d_\infty$ between adjacent windows)
involve:
\begin{enumerate}[leftmargin=1.5em]
  \item Subtracting two window vectors (768 float32 subtractions).
  \item Computing a sum of squares (768 multiplications and additions).
  \item Taking a square root.
\end{enumerate}

A distance error bound must include subtraction, squaring, reduction,
possible underflow/overflow and square root. The length768 alone does not
prove an absolute error below0.006 or make an error irrelevant near a domain
boundary. float32 is the frozen representation; descriptive distance analysis
and conservative certificate membership have different requirements. An
observed empirical tolerance is not a universal containment proof.

\subsection{Case 2: Normalization Parameter Computation}
The current train-fit normalization is stored with the frozen K=12 reference;
its preprocessing ID is
\code{aa8e1cc20829307187e0a78200f9957f7dd068b1dc26063c2720ad832bb2fd6e}.
Use the stored values and actual float32 cast/operations exactly. Historical
pure-K=6 counts such as5,693 training windows or illustrative scale0.413195
do not define current constants. Aggregating in float64 or using stable
variance algorithms can improve a new estimator, but changing the frozen
estimator/constants changes the reference and is outside this task.

The real affine formula $N(X)=(X-\mu)/\sigma$ and the executable sequence
$\widehat N(X)=\mathrm{fl}(\mathrm{fl}(X-\mu)/\sigma)$ are different objects.
For a raw-domain certificate, a justified enclosure of$\widehat N(X)$ must
fit the verified normalized domain. A normalized-input contract may instead
treat the exact stored normalized float32 tensor as the input, but must state
that boundary explicitly and preserve the real-input pipeline contract.

\subsection{Case 3: Neural Network Inference in float32}
The actual graph is padded Conv--BN(eval)--ReLU--MaxPool--Conv--BN--ReLU--
MaxPool--Conv--BN--ReLU--Flatten--Linear--ReLU--Linear12. There is no GAP head.
A bound based on $O(\sqrt{mn}\epsilon)$ or a universal$10^{-6}$ logit slack
is not a worst-case error proof for it. ReLU and max are nonexpansive for
coordinatewise error enclosures, whereas affine operators can amplify error;
BN requires frozen running statistics and its scale/epsilon evaluation.

\begin{remark}[Exact stored weights versus arithmetic error]
A finite stored float32 weight has an exact dyadic real value. Promoting that
value to float64 exactly preserves the stored weight; no unknown weight error
is introduced merely because training used float32. Uncertainty relative to
an ideal pre-rounding trained weight is a different question. The numerical
bridge needed here bounds executable operations using the stored parameters,
normalization and buffers. float64 evaluation alone does not enclose float32
rounding, and one outward rounding at the final bound does not cover errors
accumulated throughout a neural computation.
\end{remark}

If$\ell_{cj}$ is a sound ideal margin bound and
$|\widehat g_{cj}(X)-g_{cj}(X)|\le\eta_{cj}$ throughout the domain, then
$\ell_{cj}>\eta_{cj}$ suffices. If computed$\widehat\ell$ has a justified
bound error$\delta$ with$\widehat\ell-\delta\le\inf g$, require
$\widehat\ell>\delta+\eta$. All eleven competitor errors must be covered.
If these quantities are unavailable, label executable soundness CONDITIONAL
or NOT ESTABLISHED and disable formal runtime acceptance.

\subsection{Case 4: Quantization (int8, fixed-point)}

Edge deployment may quantize the network to int8 (8-bit integer) or fixed-point
arithmetic to reduce latency and energy. This changes the network function from
$f_\theta^{\text{fp32}}$ to $f_\theta^{\text{int8}}$, with differences depending on scales, zero-points, saturation, rounding and operators; no universal0.5\% error budget is assumed. Certificates computed for
$f_\theta^{\text{fp32}}$ do \textbf{not} apply to $f_\theta^{\text{int8}}$.
A separate verifier run on the quantized network is required.

% ============================================================
\section{Float32 Precision: What Is Safe and What Is Not}

\begin{center}
\small
\begin{tabularx}{\linewidth}{@{}lYY@{}}
\toprule
Operation & Safe in float32? & Notes \\
\midrule
Window tensor storage & Frozen format & Reconstruct values and identity; no universal sensor-LSB claim. \\
Normalization at inference & Needs enclosure & Match frozen values/operations; account for boundary rounding. \\
Euclidean distance & Diagnostic or bounded & No blanket error guarantee from dimension alone. \\
Mean/std estimation & Frozen estimator & Stable algorithms help new fits; do not refit this reference. \\
Softmax/log-sum-exp & Separate semantics & Subtract max improves stability; class gate uses logits. \\
Stored network weights & Exact dyadic values & Promoting preserves values; arithmetic errors remain. \\
Certification bounds & Needs validated errors & Correct operators, rounding and all-eleven margins. \\
Quantized inference & Different reference & Separate model identity and proof. \\
\bottomrule
\end{tabularx}
\end{center}

% ============================================================
\section{Summary}

The data audit confirmed (Evidence D01) that the window tensor is an exact float32
cast of the raw text values. This has three direct consequences for ECDR:
\begin{enumerate}[leftmargin=1.5em]
  \item \textbf{Reproducibility:} The reconstruction reproduces checked values under the stated procedure; hash/source receipts establish artifact identity and builder provenance.
  \item \textbf{Distance calculations:} Descriptive float32 statistics may be reported with their scope; membership near a certificate boundary requires conservative containment, not a generic adequacy claim.
  \item \textbf{Certification soundness:} The formal verifier must justify a numerical bridge to the exact executable reference; it may use another arithmetic mode if its enclosure proof is valid. Certification of a float32
        network requires either (a) interval arithmetic that bounds float32 rounding
        errors, or (b) sufficient margin slack to absorb the worst-case rounding
        error budget.
\end{enumerate}

\begin{thebibliography}{9}
\bibitem{main}
  ECDR Research Workspace. \textit{ECDR Data Analysis} (English edition).
  \code{ECDR_Data_Analysis.tex}, October~2026.
\bibitem{higham2002}
  N.J. Higham. \textit{Accuracy and Stability of Numerical Algorithms}, 2nd ed.
  SIAM, 2002.
\bibitem{ieee754}
  IEEE Standard for Floating-Point Arithmetic. \textit{IEEE Std 754-2019}.
  IEEE, 2019.
\end{thebibliography}

\end{document}
